\section{预文本任务：从图像变换中学习}

\subsection{预文本任务与下游任务}

\begin{figure}[H]
\centering
\includegraphics[width=0.95\textwidth]{slides-images/slide-009.jpg}
\caption{Self-Supervised Learning 的基本框架\protect\footnotemark}
\end{figure}
\footnotetext{Slides 第9页。}

预文本任务的标签不是人工标注，而是从数据本身自动构造出来的。常见模式是：
\begin{itemize}
    \item 输入：原始图片或视频片段；
    \item 通过某种变换或遮挡构造任务；
    \item 输出：旋转角度、patch 位置、缺失像素、颜色通道等；
    \item encoder 负责学出表征，下游任务再接一个线性层或小型头部。
\end{itemize}

\begin{knowledgebox}{评价自监督方法时看什么}
\begin{itemize}
    \item 预文本任务能否被很好地解决。
    \item 学到的表示是否有结构、是否容易聚类或可视化。
    \item 迁移到分类、检测、分割等下游任务时是否有效。
    \item 是否足够鲁棒、泛化，并且训练和推理开销合理。
\end{itemize}
\end{knowledgebox}

\subsection{旋转预测}

最经典的预文本任务之一是 \textbf{rotation prediction}。给定一张图像，将它旋转 $0^\circ, 90^\circ, 180^\circ, 270^\circ$ 之一，模型只需要预测旋转类别。

\begin{figure}[H]
\centering
\includegraphics[width=0.95\textwidth]{slides-images/slide-021.jpg}
\caption{旋转预测：四分类任务要求模型理解“视觉常识”\protect\footnotemark}
\end{figure}
\footnotetext{Slides 第21页。}

这个任务看似简单，但如果模型真的能分辨物体的朝向，说明它已经学到了一些关于“物体应该长什么样”的语义信息，而不是只看局部纹理。

\begin{importantbox}{旋转预测的直觉}
如果模型连“这张图是不是倒着的”都能分出来，它通常已经在内部形成了较好的物体级表示。
\end{importantbox}

\subsection{相对位置与拼图}

另外一类预文本任务是预测 patch 的相对位置，或者把图像切成拼图后预测正确排列。

\begin{figure}[H]
\centering
\includegraphics[width=0.95\textwidth]{slides-images/slide-025.jpg}
\caption{预测 patch 的相对位置：一种八分类的预文本任务\protect\footnotemark}
\end{figure}
\footnotetext{Slides 第25页。}

相对位置任务会迫使模型理解局部 patch 与整体结构之间的关系。拼图任务更难，因为输出空间更大，但它也更强调全局布局和对象结构。

\begin{knowledgebox}{为什么拼图任务常常要限制候选排列}
如果直接让 9 个 patch 做全排列，输出空间是 $9!$，太大了。因此很多工作会预先挑选一小组“足够不同”的排列，把问题变成有限类别分类。
\end{knowledgebox}

\subsection{修补缺失区域}

inpainting 任务会随机遮掉图像的一部分，让模型根据上下文补全缺失内容。

\begin{figure}[H]
\centering
\includegraphics[width=0.95\textwidth]{slides-images/slide-028.jpg}
\caption{Inpainting：根据上下文重建被遮挡的图像区域\protect\footnotemark}
\end{figure}
\footnotetext{Slides 第28页。}

这类方法通常是一个 encoder-decoder 结构。encoder 编码可见区域，decoder 尝试重建被 mask 的部分。训练时直接对重建误差做回归即可。

\begin{warningbox}{重建任务的一个问题}
像素级重建往往会鼓励模型输出“模糊的平均答案”，因为真实的缺失区域常常有多种合理补全。后面我们会看到，MAE 仍然沿用遮挡重建思路，但做法更现代。
\end{warningbox}

\subsection{颜色预测与 Split-Brain}

另一个经典预文本任务是 \textbf{image colorization}。做法是把图像转到 LAB 颜色空间，只输入亮度通道 $L$，让模型预测颜色通道 $A, B$。

\begin{figure}[H]
\centering
\includegraphics[width=0.95\textwidth]{slides-images/slide-034.jpg}
\caption{颜色化任务：用亮度通道预测颜色通道\protect\footnotemark}
\end{figure}
\footnotetext{Slides 第34页。}

这类任务之所以有用，是因为预测颜色不只是局部回归，它要求模型理解物体语义。比如“草通常偏绿”、“天空通常偏蓝”，这些都是物体层面的信息。

\begin{importantbox}{Split-Brain Autoencoder 的思想}
可以把输入拆成两路，让一部分通道预测另一部分通道，或者反过来。比如 RGB-D 传感器里，可以用 RGB 预测 Depth，也可以用 Depth 预测 RGB。
\end{importantbox}

\subsection{视频颜色化与跟踪}

静态图像颜色化还可以扩展到视频。给定一个有色参考帧，再去给后续灰度帧上色，模型就需要建立跨帧对应关系。

\begin{figure}[H]
\centering
\includegraphics[width=0.95\textwidth]{slides-images/slide-042.jpg}
\caption{视频颜色化：利用时间一致性把颜色从参考帧传播到后续帧\protect\footnotemark}
\end{figure}
\footnotetext{Slides 第42页。}

这里隐含的难点其实是 \textbf{tracking}。如果模型能把某个对象在不同帧中的对应区域对齐，就能把颜色或其他信息传播过去。于是，颜色化不仅是一个重建任务，也能间接逼出跟踪能力。

\begin{knowledgebox}{颜色化为何会诱发跟踪}
要把参考帧中的颜色传到目标帧，模型必须先知道哪些像素属于同一对象或同一部位。这种对应关系本质上就是跟踪。
\end{knowledgebox}

\subsection{Masked Autoencoder}

课程后半段特别强调了 \textbf{MAE}。它仍然是重建式任务，但 masking 更激进、结构更大，也更适合现代 ViT。

\begin{figure}[H]
\centering
\includegraphics[width=0.95\textwidth]{slides-images/slide-052.jpg}
\caption{Masked Autoencoders：用高比例遮挡迫使模型学出强表示\protect\footnotemark}
\end{figure}
\footnotetext{Slides 第52页。}

MAE 的关键设计是：
\begin{itemize}
    \item 先把图像切成 patch；
    \item 随机遮掉大部分 patch，常见遮挡率可高达 75\%；
    \item encoder 只处理可见 patch；
    \item decoder 用 mask token 复原完整图像；
    \item 损失只算在被 mask 的 patch 上。
\end{itemize}

\begin{importantbox}{MAE 的两个优势}
\begin{enumerate}
    \item 遮挡率高，等于对同一张图做了大量随机增强，训练信号丰富。
    \item encoder 只看少量可见 patch，因此训练效率高，表示质量也往往很好。
\end{enumerate}
\end{importantbox}

\subsection{如何评价预文本任务}

预文本任务本身并不是终点。真正重要的是它学到的 encoder 是否对下游任务有帮助。

\begin{figure}[H]
\centering
\includegraphics[width=0.95\textwidth]{slides-images/slide-067.jpg}
\caption{课程中对自监督学习的总结性提示：看下游任务的效果，而不是只看预文本损失\protect\footnotemark}
\end{figure}
\footnotetext{Slides 第67页。}

通常会看三类指标：
\begin{itemize}
    \item \textbf{linear probing}：冻结 encoder，只训练一个线性分类器；
    \item \textbf{fine-tuning}：在下游任务上继续微调部分或全部参数；
    \item \textbf{feature quality}：做可视化、聚类、检索等分析。
\end{itemize}

\begin{warningbox}{预文本任务的局限}
很多预文本任务都很“手工”，每设计一种任务就要额外调一套方法。更严重的是，学到的表示可能过度绑定某个具体任务，不够通用。
\end{warningbox}

%% ========== Contrastive ==========

